% assets/w03-euler.svg — 오일러각 셋. 무엇이 어느 축을 도는가
%
%   bash _tools/tikz2svg.sh assets/src/w03-euler.tex assets/w03-euler.svg
%
% 요점
%   오일러각은 "축 하나를 골라 그 축으로 돌리기" 를 세 번 한 것이다.
%   선박에서는 x_b(선수) 둘레가 롤, y_b(우현) 둘레가 피치, z_b(아래) 둘레가 요.
%
% 투영 — 앞왼쪽 위에서 내려다본 3/4 시점. 오른손 좌표계 (x 앞, y 우현, z 아래)
%   x_b -> (-0.75,  0.38)      y_b -> ( 0.90,  0.30)      z_b -> ( 0.00, -0.95)

\documentclass[border=6pt]{standalone}
% capstone-style.tex — 이 볼트의 TikZ 그림이 공유하는 스타일
%
% 쓰는 법:  % capstone-style.tex — 이 볼트의 TikZ 그림이 공유하는 스타일
%
% 쓰는 법:  % capstone-style.tex — 이 볼트의 TikZ 그림이 공유하는 스타일
%
% 쓰는 법:  \input{capstone-style}   (assets/src/ 안에서)
% 빌드:     bash _tools/tikz2svg.sh assets/src/w02-node-topic.tex assets/w02-node-topic.svg
%
% 왜 TikZ 인가
%   손으로 SVG 좌표를 쓰면 화살촉·선 굵기·글자 크기가 그림마다 미묘하게 달라진다.
%   그것이 "자료가 어설퍼 보인다"의 정체다. 스타일을 한 번만 정하고 상대좌표로
%   배치하면 좌표를 고쳐도 그림이 무너지지 않는다.
%   수식은 본문과 같은 TeX 글꼴로 나온다.
%
% 한글
%   ko.TeX(DVI 경로)를 쓴다. XeLaTeX 은 TikZ 그래픽을 PDF special 로 내보내는데
%   dvisvgm 이 그것을 받으려면 Ghostscript < 10.01 이 필요하다. 이 PC 에는
%   10.07 이 깔려 있어 그 경로가 막힌다. latex -> DVI -> dvisvgm 이 유일하게
%   동작하는 조합이다.

\usepackage[hangul]{kotex}
\usepackage{tikz}
\usepackage{amsmath}
\usepackage{xcolor}
\usetikzlibrary{arrows.meta, positioning, calc, fit, backgrounds, decorations.pathmorphing}

% ---- 색 --------------------------------------------------------------------
% 강조는 하나만. 나머지는 검정. 흑백 인쇄에서도 읽혀야 한다.
\definecolor{ink}{HTML}{1A1A1A}
\definecolor{soft}{HTML}{666666}
\definecolor{accent}{HTML}{7C3AED}
\definecolor{warn}{HTML}{B91C1C}

% 계층 색 — Simulink 모델의 블록 색과 같은 규칙을 쓴다
\definecolor{guid}{HTML}{1D4ED8}    % 유도
\definecolor{ctrl}{HTML}{EA580C}    % 제어
\definecolor{thr}{HTML}{D97706}     % 추진기
\definecolor{plant}{HTML}{16A34A}   % 운동모델
\definecolor{nav}{HTML}{7C3AED}     % 항법 · 신호처리
\definecolor{pathc}{HTML}{0D9488}   % 계획 경로 · 항적
\definecolor{errc}{HTML}{C2410C}    % 오차
\definecolor{flow}{HTML}{0369A1}    % 조류 · 바람

% ---- 선과 화살촉 -----------------------------------------------------------
\tikzset{
  sig/.style   = {ink, line width=0.5pt, -{Stealth[length=4pt, width=3pt]}},
  wire/.style  = {ink, line width=0.5pt},
  fb/.style    = {ink, line width=0.5pt, densely dashed,
                  -{Stealth[length=4pt, width=3pt]}},
  asig/.style  = {accent, line width=0.5pt, densely dashed,
                  -{Stealth[length=4pt, width=3pt]}},
  % 블록 — 흰 바탕, 직각, 검은 테두리. 색은 테두리에만 준다
  blk/.style   = {draw=ink, line width=0.55pt, fill=white,
                  minimum width=18mm, minimum height=8.5mm, inner sep=2pt, font=\small},
  gblk/.style  = {blk, draw=guid},
  cblk/.style  = {blk, draw=ctrl},
  tblk/.style  = {blk, draw=thr},
  pblk/.style  = {blk, draw=plant},
  nblk/.style  = {blk, draw=nav},
  % 합산점 · 분기점
  sum/.style   = {draw=ink, line width=0.55pt, fill=white, circle,
                  inner sep=0pt, minimum size=4.6mm, font=\scriptsize},
  dot/.style   = {ink, fill=ink, circle, inner sep=0pt, minimum size=1.5mm},
  % 글자
  sname/.style = {font=\small, inner sep=1.5pt},
  note/.style  = {font=\scriptsize, soft, inner sep=1.5pt},
  wnote/.style = {font=\scriptsize, warn, inner sep=1.5pt},
  anote/.style = {font=\scriptsize, accent, inner sep=1.5pt},
  axis/.style  = {ink, line width=0.5pt},
  tick/.style  = {font=\scriptsize, soft},
  % 묶음 상자 — 서브시스템 표시
  grp/.style   = {draw=soft, line width=0.4pt, densely dotted, rounded corners=1mm,
                  inner sep=3mm},
}
   (assets/src/ 안에서)
% 빌드:     bash _tools/tikz2svg.sh assets/src/w02-node-topic.tex assets/w02-node-topic.svg
%
% 왜 TikZ 인가
%   손으로 SVG 좌표를 쓰면 화살촉·선 굵기·글자 크기가 그림마다 미묘하게 달라진다.
%   그것이 "자료가 어설퍼 보인다"의 정체다. 스타일을 한 번만 정하고 상대좌표로
%   배치하면 좌표를 고쳐도 그림이 무너지지 않는다.
%   수식은 본문과 같은 TeX 글꼴로 나온다.
%
% 한글
%   ko.TeX(DVI 경로)를 쓴다. XeLaTeX 은 TikZ 그래픽을 PDF special 로 내보내는데
%   dvisvgm 이 그것을 받으려면 Ghostscript < 10.01 이 필요하다. 이 PC 에는
%   10.07 이 깔려 있어 그 경로가 막힌다. latex -> DVI -> dvisvgm 이 유일하게
%   동작하는 조합이다.

\usepackage[hangul]{kotex}
\usepackage{tikz}
\usepackage{amsmath}
\usepackage{xcolor}
\usetikzlibrary{arrows.meta, positioning, calc, fit, backgrounds, decorations.pathmorphing}

% ---- 색 --------------------------------------------------------------------
% 강조는 하나만. 나머지는 검정. 흑백 인쇄에서도 읽혀야 한다.
\definecolor{ink}{HTML}{1A1A1A}
\definecolor{soft}{HTML}{666666}
\definecolor{accent}{HTML}{7C3AED}
\definecolor{warn}{HTML}{B91C1C}

% 계층 색 — Simulink 모델의 블록 색과 같은 규칙을 쓴다
\definecolor{guid}{HTML}{1D4ED8}    % 유도
\definecolor{ctrl}{HTML}{EA580C}    % 제어
\definecolor{thr}{HTML}{D97706}     % 추진기
\definecolor{plant}{HTML}{16A34A}   % 운동모델
\definecolor{nav}{HTML}{7C3AED}     % 항법 · 신호처리
\definecolor{pathc}{HTML}{0D9488}   % 계획 경로 · 항적
\definecolor{errc}{HTML}{C2410C}    % 오차
\definecolor{flow}{HTML}{0369A1}    % 조류 · 바람

% ---- 선과 화살촉 -----------------------------------------------------------
\tikzset{
  sig/.style   = {ink, line width=0.5pt, -{Stealth[length=4pt, width=3pt]}},
  wire/.style  = {ink, line width=0.5pt},
  fb/.style    = {ink, line width=0.5pt, densely dashed,
                  -{Stealth[length=4pt, width=3pt]}},
  asig/.style  = {accent, line width=0.5pt, densely dashed,
                  -{Stealth[length=4pt, width=3pt]}},
  % 블록 — 흰 바탕, 직각, 검은 테두리. 색은 테두리에만 준다
  blk/.style   = {draw=ink, line width=0.55pt, fill=white,
                  minimum width=18mm, minimum height=8.5mm, inner sep=2pt, font=\small},
  gblk/.style  = {blk, draw=guid},
  cblk/.style  = {blk, draw=ctrl},
  tblk/.style  = {blk, draw=thr},
  pblk/.style  = {blk, draw=plant},
  nblk/.style  = {blk, draw=nav},
  % 합산점 · 분기점
  sum/.style   = {draw=ink, line width=0.55pt, fill=white, circle,
                  inner sep=0pt, minimum size=4.6mm, font=\scriptsize},
  dot/.style   = {ink, fill=ink, circle, inner sep=0pt, minimum size=1.5mm},
  % 글자
  sname/.style = {font=\small, inner sep=1.5pt},
  note/.style  = {font=\scriptsize, soft, inner sep=1.5pt},
  wnote/.style = {font=\scriptsize, warn, inner sep=1.5pt},
  anote/.style = {font=\scriptsize, accent, inner sep=1.5pt},
  axis/.style  = {ink, line width=0.5pt},
  tick/.style  = {font=\scriptsize, soft},
  % 묶음 상자 — 서브시스템 표시
  grp/.style   = {draw=soft, line width=0.4pt, densely dotted, rounded corners=1mm,
                  inner sep=3mm},
}
   (assets/src/ 안에서)
% 빌드:     bash _tools/tikz2svg.sh assets/src/w02-node-topic.tex assets/w02-node-topic.svg
%
% 왜 TikZ 인가
%   손으로 SVG 좌표를 쓰면 화살촉·선 굵기·글자 크기가 그림마다 미묘하게 달라진다.
%   그것이 "자료가 어설퍼 보인다"의 정체다. 스타일을 한 번만 정하고 상대좌표로
%   배치하면 좌표를 고쳐도 그림이 무너지지 않는다.
%   수식은 본문과 같은 TeX 글꼴로 나온다.
%
% 한글
%   ko.TeX(DVI 경로)를 쓴다. XeLaTeX 은 TikZ 그래픽을 PDF special 로 내보내는데
%   dvisvgm 이 그것을 받으려면 Ghostscript < 10.01 이 필요하다. 이 PC 에는
%   10.07 이 깔려 있어 그 경로가 막힌다. latex -> DVI -> dvisvgm 이 유일하게
%   동작하는 조합이다.

\usepackage[hangul]{kotex}
\usepackage{tikz}
\usepackage{amsmath}
\usepackage{xcolor}
\usetikzlibrary{arrows.meta, positioning, calc, fit, backgrounds, decorations.pathmorphing}

% ---- 색 --------------------------------------------------------------------
% 강조는 하나만. 나머지는 검정. 흑백 인쇄에서도 읽혀야 한다.
\definecolor{ink}{HTML}{1A1A1A}
\definecolor{soft}{HTML}{666666}
\definecolor{accent}{HTML}{7C3AED}
\definecolor{warn}{HTML}{B91C1C}

% 계층 색 — Simulink 모델의 블록 색과 같은 규칙을 쓴다
\definecolor{guid}{HTML}{1D4ED8}    % 유도
\definecolor{ctrl}{HTML}{EA580C}    % 제어
\definecolor{thr}{HTML}{D97706}     % 추진기
\definecolor{plant}{HTML}{16A34A}   % 운동모델
\definecolor{nav}{HTML}{7C3AED}     % 항법 · 신호처리
\definecolor{pathc}{HTML}{0D9488}   % 계획 경로 · 항적
\definecolor{errc}{HTML}{C2410C}    % 오차
\definecolor{flow}{HTML}{0369A1}    % 조류 · 바람

% ---- 선과 화살촉 -----------------------------------------------------------
\tikzset{
  sig/.style   = {ink, line width=0.5pt, -{Stealth[length=4pt, width=3pt]}},
  wire/.style  = {ink, line width=0.5pt},
  fb/.style    = {ink, line width=0.5pt, densely dashed,
                  -{Stealth[length=4pt, width=3pt]}},
  asig/.style  = {accent, line width=0.5pt, densely dashed,
                  -{Stealth[length=4pt, width=3pt]}},
  % 블록 — 흰 바탕, 직각, 검은 테두리. 색은 테두리에만 준다
  blk/.style   = {draw=ink, line width=0.55pt, fill=white,
                  minimum width=18mm, minimum height=8.5mm, inner sep=2pt, font=\small},
  gblk/.style  = {blk, draw=guid},
  cblk/.style  = {blk, draw=ctrl},
  tblk/.style  = {blk, draw=thr},
  pblk/.style  = {blk, draw=plant},
  nblk/.style  = {blk, draw=nav},
  % 합산점 · 분기점
  sum/.style   = {draw=ink, line width=0.55pt, fill=white, circle,
                  inner sep=0pt, minimum size=4.6mm, font=\scriptsize},
  dot/.style   = {ink, fill=ink, circle, inner sep=0pt, minimum size=1.5mm},
  % 글자
  sname/.style = {font=\small, inner sep=1.5pt},
  note/.style  = {font=\scriptsize, soft, inner sep=1.5pt},
  wnote/.style = {font=\scriptsize, warn, inner sep=1.5pt},
  anote/.style = {font=\scriptsize, accent, inner sep=1.5pt},
  axis/.style  = {ink, line width=0.5pt},
  tick/.style  = {font=\scriptsize, soft},
  % 묶음 상자 — 서브시스템 표시
  grp/.style   = {draw=soft, line width=0.4pt, densely dotted, rounded corners=1mm,
                  inner sep=3mm},
}


\begin{document}
\begin{tikzpicture}[x=1mm, y=1mm]

\pgfmathsetmacro{\S}{1.7}                    % 1 단위 = 1.7 mm
% 투영 — #1 = x(앞), #2 = y(우현), #3 = z(아래)
\newcommand{\PX}[3]{\S*(-0.75*#1 + 0.90*#2 + 0.00*#3)}
\newcommand{\PY}[3]{\S*( 0.38*#1 + 0.30*#2 - 0.95*#3)}

% ---------------------------------------------------------------------------
% \hull  선체 — 갑판 오각형. 선수가 +x
% ---------------------------------------------------------------------------
\newcommand{\hull}{%
  \draw[ink, line width=0.5pt, fill=white, fill opacity=0.9]
       ({\PX{9}{0}{0}},{\PY{9}{0}{0}}) --
       ({\PX{4}{4}{0}},{\PY{4}{4}{0}}) --
       ({\PX{-6}{4}{0}},{\PY{-6}{4}{0}}) --
       ({\PX{-6}{-4}{0}},{\PY{-6}{-4}{0}}) --
       ({\PX{4}{-4}{0}},{\PY{4}{-4}{0}}) -- cycle;}

% ---------------------------------------------------------------------------
% \axes  몸체 3축
% ---------------------------------------------------------------------------
\newcommand{\axes}[1]{%
  \draw[#1, line width=0.45pt, -{Stealth[length=3.5pt,width=2.5pt]}]
       (0,0) -- ({\PX{14}{0}{0}},{\PY{14}{0}{0}});
  \node[font=\scriptsize, #1] at ({\PX{17}{0}{0}},{\PY{17}{0}{0}}) {$x_b$};
  \draw[#1, line width=0.45pt, -{Stealth[length=3.5pt,width=2.5pt]}]
       (0,0) -- ({\PX{0}{11}{0}},{\PY{0}{11}{0}});
  \node[font=\scriptsize, #1] at ({\PX{0}{14}{0}},{\PY{0}{14}{0}}) {$y_b$};
  \draw[#1, line width=0.45pt, -{Stealth[length=3.5pt,width=2.5pt]}]
       (0,0) -- ({\PX{0}{0}{10}},{\PY{0}{0}{10}});
  \node[font=\scriptsize, #1] at ({\PX{0}{0}{13}},{\PY{0}{0}{13}}) {$z_b$};}

% ---------------------------------------------------------------------------
% \panel{x}{y}{제목}{각 기호}{설명}{강조축 0=x 1=y 2=z}
% ---------------------------------------------------------------------------
\newcommand{\panel}[6]{%
\begin{scope}[shift={(#1,#2)}]
  \hull
  \axes{soft}
  % 도는 축만 진하게 다시 긋는다
  \ifnum#6=0
    \draw[ctrl, line width=1.1pt, -{Stealth[length=4pt,width=3pt]}]
         (0,0) -- ({\PX{14}{0}{0}},{\PY{14}{0}{0}});
    \draw[ctrl, line width=0.8pt, -{Stealth[length=3.5pt,width=2.5pt]}]
         ({\PX{7}{5}{0}},{\PY{7}{5}{0}})
         .. controls ({\PX{7}{5}{-4}},{\PY{7}{5}{-4}}) and ({\PX{7}{-2}{-6}},{\PY{7}{-2}{-6}})
         .. ({\PX{7}{-5}{-2}},{\PY{7}{-5}{-2}});
  \fi
  \ifnum#6=1
    \draw[ctrl, line width=1.1pt, -{Stealth[length=4pt,width=3pt]}]
         (0,0) -- ({\PX{0}{11}{0}},{\PY{0}{11}{0}});
    \draw[ctrl, line width=0.8pt, -{Stealth[length=3.5pt,width=2.5pt]}]
         ({\PX{7}{5}{0}},{\PY{7}{5}{0}})
         .. controls ({\PX{9}{5}{-4}},{\PY{9}{5}{-4}}) and ({\PX{2}{5}{-7}},{\PY{2}{5}{-7}})
         .. ({\PX{-4}{5}{-4}},{\PY{-4}{5}{-4}});
  \fi
  \ifnum#6=2
    \draw[ctrl, line width=1.1pt, -{Stealth[length=4pt,width=3pt]}]
         (0,0) -- ({\PX{0}{0}{10}},{\PY{0}{0}{10}});
    \draw[ctrl, line width=0.8pt, -{Stealth[length=3.5pt,width=2.5pt]}]
         ({\PX{10}{4}{0}},{\PY{10}{4}{0}})
         .. controls ({\PX{6}{9}{0}},{\PY{6}{9}{0}}) and ({\PX{-2}{10}{0}},{\PY{-2}{10}{0}})
         .. ({\PX{-7}{7}{0}},{\PY{-7}{7}{0}});
  \fi
  \node[font=\small\bfseries, anchor=west] at (-20,26) {#3};
  \node[font=\small, ctrl, anchor=west] at (-20,20.5) {#4};
  \node[note, anchor=west, text width=44mm] at (-20,-19) {#5};
\end{scope}}

\panel{0}{0}{롤 (Roll)}{$\phi$ — $x_b$ 둘레}
  {좌우로 기울어짐. 파도를 옆에서 맞으면 커진다}{0}

\panel{62}{0}{피치 (Pitch)}{$\theta$ — $y_b$ 둘레}
  {앞뒤로 기울어짐. 뱃머리가 들렸다 잠겼다 한다}{1}

\panel{124}{0}{요 (Yaw)}{$\psi$ — $z_b$ 둘레}
  {뱃머리 방향이 돌아감. \textbf{선수각}이 이것이다}{2}

% =========================== 정리 ==========================================
\node[font=\small\bfseries, anchor=west] at (-20,-32) {세 번 돌린다 — 순서가 곧 정의다};

\node[font=\small, anchor=west] at (-20,-42)
     {$\mathbf{R}(\phi,\theta,\psi) \;=\; \mathbf{R}_z(\psi)\,\mathbf{R}_y(\theta)\,\mathbf{R}_x(\phi)$};

\node[note, anchor=west, text width=74mm] at (56,-38)
     {오른쪽부터 읽는다 — 롤, 피치, 요 순서로 돌린다. 순서를 바꾸면 \textbf{다른 자세}가
      되므로, 오일러각은 "세 숫자" 가 아니라 "세 숫자와 순서" 다.
      항공·선박은 $z$-$y$-$x$ (요-피치-롤) 를 쓴다};

\node[wnote, anchor=west, text width=150mm] at (-20,-54)
     {피치가 $\pm 90^{\circ}$ 에 닿으면 롤 축과 요 축이 겹쳐 자유도 하나를 잃는다
      (\textbf{짐벌락}). 수상선은 피치가 몇 도를 넘지 않아 실무에서 만나지 않지만,
      저장·전송에는 이 문제가 없는 쿼터니언을 쓴다};

\end{tikzpicture}
\end{document}
